% Lecture notes: Week 12, Lecture 1 -- Statically Indeterminate Bending Problems
% To hide this lecture from the website, add a line containing only:  % draft
% To set the date shown on the website, add a line like:  % posted: 2026-09-02
\documentclass[11pt]{article}
\usepackage{mech2030notes}
\begin{document}

% ##################################################################
%  WEEK 12, LECTURE 1
% ##################################################################
\lecture{12}{1}{Statically Indeterminate Bending Problems}

\section*{Example}

\noindent
\begin{minipage}[c]{0.55\textwidth}
\centering
\begin{tikzpicture}
  % scale: 2.5 cm per m
  \lwall{0}{-0.9}{0.9}
  \draw[beam] (0,-0.15) rectangle (5,0.15);
  \node[left] at (-0.4,0) {$A$};
  % roller at B
  \draw[thick, fill=white] (2.5,-0.15) -- (2.2,-0.55) -- (2.8,-0.55) -- cycle;
  \foreach \dx in {-0.18,0,0.18} \draw[thick] ({2.5+\dx},-0.64) circle (0.09);
  \ground{2.05}{2.95}{-0.73}
  \node[above] at (2.5,0.15) {$B$};
  \draw[force] (5,1.3) node[above] {$P = \SI{10}{kN}$} -- (5,0.17);
  \draw[dim] (0,-1.2) -- (2.5,-1.2) node[midway, fill=white] {$L = \SI{1}{m}$};
  \draw[dim] (2.5,-1.2) -- (5,-1.2) node[midway, fill=white] {$L = \SI{1}{m}$};
\end{tikzpicture}
\end{minipage}%
\hfill
\begin{minipage}[c]{0.42\textwidth}
\textbf{Setup:} A beam is loaded as shown, with $P = \SI{10}{kN}$, $L = \SI{1}{m}$, and
$EI$ constant.

\smallskip
\textbf{Find:} Determine the reactions at the supports.
\end{minipage}

\subsection*{FBD}

\begin{center}
\begin{tikzpicture}
  \draw[thick] (0,0) -- (5,0);
  \node[pin] at (0,0) {}; \node[pin] at (2.5,0) {}; \node[pin] at (5,0) {};
  \draw[force] ([shift={(20:0.5)}]0,0) arc (20:160:0.5);
  \node[left] at (-0.5,0.3) {$M_A$};
  \draw[force] (0,-1.1) node[below] {$A_y$} -- (0,-0.08);
  \draw[force] (2.5,-1.1) node[below] {$B_y$} -- (2.5,-0.08);
  \draw[force] (5,1.1) node[above] {$\SI{10}{kN}$} -- (5,0.08);
\end{tikzpicture}
\end{center}

\begin{itemize}[label=$\rightarrow$]
  \item 3 unknowns: $M_A$, $A_y$, $B_y$.
  \item 2 equations: $\sum F_y = 0$, $\sum M = 0$.
\end{itemize}
\begin{center}
  $\Rightarrow$ \emph{Statically indeterminate.}
\end{center}

\noindent
$\rightarrow$ Introduce a \emph{compatibility condition} as the 3rd equation:
\[
  \boxed{v_B = 0} \qquad \text{(no deflection at the roller $B$)}
\]

% ------------------------------------------------------------------
\subsection*{Superposition}
% ------------------------------------------------------------------

\noindent
Remove the roller and replace it with the (unknown) reaction $B_y$. The deflection at $B$ is
the sum of the deflections from each load acting on the cantilever separately:

\begin{center}
\begin{tikzpicture}[baseline=(b), scale=0.8]
  \coordinate (b) at (0,0);
  \lwall{0}{-0.9}{0.9}
  \draw[beam] (0,-0.15) rectangle (5,0.15);
  \node[pin] at (2.5,0) {}; \node[above] at (2.5,0.15) {$B$};
  \draw[force] (2.5,-1.2) node[below] {$B_y$} -- (2.5,-0.17);
  \node at (4,-0.9) {\circnum{I}};
  \draw[dim] (0,-1.95) -- (2.5,-1.95) node[midway, fill=white] {$\SI{1}{m}$};
\end{tikzpicture}
\quad{\Large $+$}\quad
\begin{tikzpicture}[baseline=(b), scale=0.8]
  \coordinate (b) at (0,0);
  \lwall{0}{-0.9}{0.9}
  \draw[beam] (0,-0.15) rectangle (5,0.15);
  \node[pin] at (2.5,0) {}; \node[above] at (2.5,0.15) {$B$};
  \draw[force] (5,1.2) node[above] {$\SI{10}{kN}$} -- (5,0.17);
  \node at (1.25,-0.9) {\circnum{II}};
  \draw[dim] (0,-1.95) -- (2.5,-1.95) node[midway, fill=white] {$\SI{1}{m}$};
  \draw[dim] (2.5,-1.95) -- (5,-1.95) node[midway, fill=white] {$\SI{1}{m}$};
\end{tikzpicture}
\end{center}

\noindent
From the cantilever table (derived in Week~7, Lecture~3), for a tip load $P$ on a
cantilever of length $L$:
\[
  v_{\max} = -\frac{PL^3}{3EI},
  \qquad\qquad
  v(x) = -\frac{Px^2}{6EI}\,(3L - x) \quad \text{(elastic curve)}
\]
\begin{itemize}
  \item \circnum{I}: $B_y$ acts \emph{upward} at $x = \SI{1}{m}$. The beam beyond $B$ is
        unloaded, so $v_B$ is the tip deflection of a $\SI{1}{m}$ cantilever (positive,
        since $B_y$ points up).
  \item \circnum{II}: the $\SI{10}{kN}$ load acts at the tip of the full $\SI{2}{m}$
        cantilever, so use the elastic curve with $L = \SI{2}{m}$, evaluated at
        $x = \SI{1}{m}$.
\end{itemize}

\[
  v_B = 0 = \frac{B_y(\SI{1}{m})^3}{3EI}
  + \left(-\frac{(\SI{10}{kN})(\SI{1}{m})^2}{6EI}\bigl(3(\SI{2}{m}) - (\SI{1}{m})\bigr)\right)
\]
Cancelling $EI$:
\[
  v_B = 0 = \frac{B_y(\SI{1}{m})^3}{3} - \frac{(\SI{10}{kN})(\SI{1}{m})^2}{6}(\SI{5}{m})
\]
\[
  \frac{B_y(\SI{1}{m})^3}{3} = \frac{(\SI{10}{kN})(\SI{1}{m})^2(\SI{5}{m})}{6}
\]
\[
  B_y = \left(\frac{3}{6}\right)\frac{(\SI{10}{kN})(\SI{1}{m})^2(\SI{5}{m})}{(\SI{1}{m})^3}
      = \left(\frac12\right)\SI{50}{kN}
  \qquad\Rightarrow\qquad
  \fbox{\fbox{$B_y = \SI{25}{kN}$}}
\]

% ------------------------------------------------------------------
\subsection*{Global FBD}
% ------------------------------------------------------------------

\[
  \textstyle\sum F_y = 0 \quad\rightarrow\quad A_y + B_y - \SI{10}{kN} = 0
\]
\[
  A_y = \SI{10}{kN} - B_y = \SI{10}{kN} - \SI{25}{kN}
  \qquad\Rightarrow\qquad
  \fbox{\fbox{$A_y = -\SI{15}{kN}$}}
\]

\[
  \textstyle\sum M^{(A)} = 0 \quad\rightarrow\quad
  M_A + (\SI{1}{m})B_y - (\SI{2}{m})(\SI{10}{kN}) = 0
\]
\[
  M_A = \SI{20}{kN.m} - B_y(\SI{1}{m}) = \SI{20}{kN.m} - \SI{25}{kN.m}
  \qquad\Rightarrow\qquad
  \fbox{\fbox{$M_A = -\SI{5}{kN.m}$}}
\]

\noindent
\textbf{Final FBD} (all forces drawn in their actual directions):

\begin{center}
\begin{tikzpicture}
  \draw[thick] (0,0) -- (5,0);
  \node[pin] at (0,0) {}; \node[pin] at (2.5,0) {}; \node[pin] at (5,0) {};
  % M_A = -5 kN.m: actually clockwise
  \draw[force] ([shift={(160:0.5)}]0,0) arc (160:20:0.5);
  \node[left] at (-0.55,0.35) {$\SI{5}{kN.m}$};
  \draw[force] (0,-0.08) -- (0,-1.1) node[below] {$\SI{15}{kN}$};
  \draw[force] (2.5,-1.1) node[below] {$\SI{25}{kN}$} -- (2.5,-0.08);
  \draw[force] (5,1.1) node[above] {$\SI{10}{kN}$} -- (5,0.08);
\end{tikzpicture}
\end{center}

\end{document}
